% =============================================================================
% Supplemental Text SX: Physical-Space Derivation of Prediction Equations
% Insert as a new Text section in Supporting Information S1, after Text S2.
% Notation matches main text: tildes for standardized, hats for kernel-integrated
% standardized features, overbars for kernel-integrated physical features.
% =============================================================================

\textbf{Text S3.  Physical-Space Derivation of Prediction Equations}

\vspace{0.5em}
\textbf{General Prediction Equation.}
All models in the hierarchy share a common prediction structure.  Each model
learns a function $f(\mathbf{x})$ of standardized predictors and produces the
standardized log-precipitation prediction
%
\begin{equation}
  z \;=\; z_{\min} + \max\!\bigl(f(\mathbf{x}),\;0\bigr),
  \qquad
  z_{\min} = -\,\frac{\mu_y}{s_y}\,,
  \label{eq:s-zfloor}
\end{equation}
%
where $\mu_y$ and $s_y$ are the training-set mean and standard deviation of
$y = \log(1 + P)$.  Physical precipitation is recovered by inverting the
standardization and log-transformation:
%
\begin{align}
  P &= \exp\!\bigl(\mu_y + s_y\,z\bigr) - 1 \notag\\[4pt]
    &= \exp\!\Bigl(\mu_y + s_y\bigl[z_{\min}
       + \max\!\bigl(f(\mathbf{x}),\;0\bigr)\bigr]\Bigr) - 1 \notag\\[4pt]
    &= \exp\!\Bigl(\mu_y + s_y\Bigl[-\frac{\mu_y}{s_y}
       + \max\!\bigl(f(\mathbf{x}),\;0\bigr)\Bigr]\Bigr) - 1 \notag\\[4pt]
    &= \exp\!\Bigl(\mu_y - \mu_y
       + s_y\max\!\bigl(f(\mathbf{x}),\;0\bigr)\Bigr) - 1 \notag\\[4pt]
    &= \exp\!\Bigl(s_y\max\!\bigl(f(\mathbf{x}),\;0\bigr)\Bigr) - 1\,.
  \label{eq:s-cancel}
\end{align}
%
Because $s_y > 0$ and $\max(f(\mathbf{x}),\,0) \geq 0$, the argument of the
exponential is non-negative, so $P \geq 0$ and the outer $\max(\cdot,\,0)$ is
redundant.  When $f(\mathbf{x}) \leq 0$ the prediction collapses to $P = 0$,
and when $f(\mathbf{x}) > 0$ precipitation increases as
%
\begin{equation}
  P \;=\; \exp\!\bigl(s_y\,f(\mathbf{x})\bigr) - 1\,.
  \label{eq:s-pred-positive}
\end{equation}
%
Combining both cases yields the prediction equation (Equation~1 in the main
text):
%
\begin{equation}
  \boxed{P \;=\; \max\!\Bigl(\exp\!\bigl(s_y\,f(\mathbf{x})\bigr)
         - 1,\;\;0\Bigr).}
  \label{eq:s-pred}
\end{equation}

\vspace{0.5em}
\textbf{Standardization of Kernel-Integrated Features.}
In the main text (Section~3.3), vertically varying predictors
$\{\varphi_i(\sigma)\}$ are standardized before kernel integration, producing
kernel-integrated standardized features $\hat{\varphi}_i$ (Equation~3).  The
physical-space derivations below require expressing $\hat{\varphi}_i$ in terms
of the corresponding kernel-integrated physical feature
%
\begin{equation}
  \bar{\varphi}_i \;=\; \sum_{m=1}^{N_\sigma}
    k_i(\sigma_m)\,\varphi_i(\sigma_m)\,\Delta\sigma_m\,,
  \label{eq:s-phys-ki}
\end{equation}
%
where $k_i(\sigma_m)$ is the learned kernel weight and $\Delta\sigma_m$ the
vertical quadrature weight.

Because each predictor is standardized using a single global mean
$\mu_{\varphi_i}$ and standard deviation $s_{\varphi_i}$ computed across all
grid points and vertical levels in the training set, the standardized profile
is $\tilde{\varphi}_i(\sigma_m) = (\varphi_i(\sigma_m) -
\mu_{\varphi_i})/s_{\varphi_i}$.  Substituting into the kernel integration
(Equation~3) and using the normalization condition $\sum_m k_i(\sigma_m)\,
\Delta\sigma_m = 1$ gives
%
\begin{align}
  \hat{\varphi}_i
    &= \sum_{m} k_i(\sigma_m)\,
       \frac{\varphi_i(\sigma_m) - \mu_{\varphi_i}}{s_{\varphi_i}}
       \,\Delta\sigma_m \notag\\[4pt]
    &= \frac{1}{s_{\varphi_i}}\sum_{m} k_i(\sigma_m)\,
       \varphi_i(\sigma_m)\,\Delta\sigma_m
       \;-\; \frac{\mu_{\varphi_i}}{s_{\varphi_i}}\underbrace{%
       \sum_{m} k_i(\sigma_m)\,\Delta\sigma_m}_{=\,1} \notag\\[4pt]
    &= \frac{\bar{\varphi}_i - \mu_{\varphi_i}}{s_{\varphi_i}}\,.
  \label{eq:s-ki-standardize}
\end{align}
%
The kernel-integrated standardized feature is therefore identical in form to
ordinary scalar standardization applied to the kernel-integrated physical
feature.  This identity holds because the standardization statistics are global
scalars (not level-dependent) and the kernels are normalized to integrate to
unity.  It allows all symbolic regression equations to be converted to
physical-space forms by direct algebraic substitution without retraining.

\vspace{1.0em}
\textbf{SR-BL: Buoyancy-Only Equation.}
SR-BL operates on a single scalar predictor, the lower-tropospheric buoyancy
$B_L$, standardized as $\tilde{B}_L = (B_L - \mu_{B_L})/s_{B_L}$.  The
discovered equation in standardized space (Table~S2) is
%
\begin{equation}
  f(\mathbf{x}) \;=\; (\tilde{B}_L + a)^3 + b\,,
  \label{eq:s-srbl-std}
\end{equation}
%
with optimized constants $a$ and $b$ (Table~SX).

Substituting $\tilde{B}_L = (B_L - \mu_{B_L})/s_{B_L}$:
%
\begin{align}
  f &= \left(\frac{B_L - \mu_{B_L}}{s_{B_L}} + a\right)^{\!3} + b
       \notag\\[4pt]
    &= \left(\frac{B_L - (\mu_{B_L} - a\,s_{B_L})}{s_{B_L}}\right)^{\!3} + b
       \notag\\[4pt]
    &= \frac{1}{s_{B_L}^3}\bigl(B_L - B_c\bigr)^3 + b\,,
  \label{eq:s-srbl-expand}
\end{align}
%
where $B_c = \mu_{B_L} - a\,s_{B_L}$ is the critical buoyancy threshold.
Substituting into Equation~\ref{eq:s-pred}:
%
\begin{align}
  P &= \max\!\Bigl(\exp\!\Bigl(
       \frac{s_y}{s_{B_L}^3}(B_L - B_c)^3 + s_y\,b
       \Bigr) - 1,\;\;0\Bigr) \notag\\[4pt]
    &= \max\!\Bigl(\exp\!\bigl(
       \alpha\,(B_L - B_c)^3 + \beta\bigr) - 1,\;\;0\Bigr),
  \label{eq:s-srbl-phys}
\end{align}
%
with three interpretable physical constants:
%
\begin{equation}
  \begin{aligned}
    B_c    &= \mu_{B_L} - a\,s_{B_L}
             &&\text{(critical buoyancy threshold)}\,,\\
    \alpha &= s_y\,/\,s_{B_L}^3
             &&\text{(exponential sensitivity to buoyancy departures)}\,,\\
    \beta  &= s_y\,b
             &&\text{(log-space magnitude offset)}\,.
  \end{aligned}
  \label{eq:s-srbl-constants}
\end{equation}
%
Precipitation onset ($P > 0$) occurs when $\alpha\,(B_L - B_c)^3 + \beta > 0$,
giving a critical onset threshold $B_L > B_c - (\beta/\alpha)^{1/3}$.

\vspace{0.3em}
\noindent\textit{Numerical values.}  With the optimized constants from
Table~S2 ($a = \text{[TBD]}$, $b = \text{[TBD]}$) and the training-set
statistics $\mu_{B_L} = \text{[TBD]}$~m\,s$^{-2}$,
$s_{B_L} = \text{[TBD]}$~m\,s$^{-2}$, $s_y = \text{[TBD]}$:
%
\begin{equation*}
  B_c = \text{[TBD]~m\,s}^{-2}\,,\qquad
  \alpha = \text{[TBD]~(m\,s}^{-2})^{-3}\,,\qquad
  \beta = \text{[TBD]}\,.
\end{equation*}

% -------------------------------------------------------------------------

\vspace{1.0em}
\textbf{SR-ATM: Atmospheric Thermodynamic Equation.}
SR-ATM uses the Gaussian-kernel-integrated features $\widehat{\mathrm{RH}}$,
$\hat{\theta}_e$, and $\hat{\theta}_e^*$ as predictors.  The discovered
equation in standardized space (Table~S2) is
%
\begin{equation}
  f(\mathbf{x}) \;=\; a\,\max\!\bigl(\widehat{\mathrm{RH}},\;\;
    \hat{\theta}_e - b\,\hat{\theta}_e^* - c\bigr)^3,
  \label{eq:s-sratm-std}
\end{equation}
%
with optimized constants $a$, $b$, and $c$ (Table~SX).

Using Equation~\ref{eq:s-ki-standardize}, each kernel-integrated standardized
feature equals the standardized form of its physical counterpart:
$\widehat{\mathrm{RH}} = (\overline{\mathrm{RH}} -
\mu_{\mathrm{RH}})/s_{\mathrm{RH}}$, $\hat{\theta}_e = (\bar{\theta}_e -
\mu_{\theta_e})/s_{\theta_e}$, and $\hat{\theta}_e^* = (\bar{\theta}_e^* -
\mu_{\theta_e^*})/s_{\theta_e^*}$.  Substituting:
%
\begin{align}
  f &= a\,\max\!\Biggl(
       \frac{\overline{\mathrm{RH}} - \mu_{\mathrm{RH}}}{s_{\mathrm{RH}}},
       \;\;\;
       \frac{\bar{\theta}_e - \mu_{\theta_e}}{s_{\theta_e}}
       - b\,\frac{\bar{\theta}_e^* - \mu_{\theta_e^*}}{s_{\theta_e^*}}
       - c
       \Biggr)^{\!3}.
  \label{eq:s-sratm-expand}
\end{align}
%
Define the two arguments of the $\max$ operator in terms of physical features:
%
\begin{align}
  A_1 &= \frac{\overline{\mathrm{RH}} - \mu_{\mathrm{RH}}}
              {s_{\mathrm{RH}}}\,,
  \label{eq:s-A1}\\[6pt]
  A_2 &= \frac{\bar{\theta}_e - \mu_{\theta_e}}{s_{\theta_e}}
       - b\,\frac{\bar{\theta}_e^* - \mu_{\theta_e^*}}{s_{\theta_e^*}}
       - c
       \;=\; \frac{\bar{\theta}_e}{s_{\theta_e}}
       - \frac{b\,\bar{\theta}_e^*}{s_{\theta_e^*}}
       - \underbrace{\biggl(\frac{\mu_{\theta_e}}{s_{\theta_e}}
       - \frac{b\,\mu_{\theta_e^*}}{s_{\theta_e^*}} + c\biggr)}_{\displaystyle
       \equiv\; d_0}\,.
  \label{eq:s-A2}
\end{align}
%
The full prediction is then
%
\begin{equation}
  \boxed{P \;=\; \max\!\Bigl(\exp\!\bigl(
    s_y\,a\;\max(A_1,\,A_2)^3\bigr) - 1,\;\;0\Bigr),}
  \label{eq:s-sratm-phys}
\end{equation}
%
where $A_1$ and $A_2$ are defined in
Equations~\ref{eq:s-A1}--\ref{eq:s-A2}, $\overline{\mathrm{RH}}$,
$\bar{\theta}_e$, and $\bar{\theta}_e^*$ are the kernel-integrated physical
features (Equation~\ref{eq:s-phys-ki} evaluated with the learned Gaussian
kernels in Table~S1), and $s_y$, $\mu_{\mathrm{RH}}$, $s_{\mathrm{RH}}$,
$\mu_{\theta_e}$, $s_{\theta_e}$, $\mu_{\theta_e^*}$, $s_{\theta_e^*}$ are
training-set statistics.

The $\max(A_1,\,A_2)$ structure selects, at each grid point and timestep,
whichever of two thermodynamic pathways is more anomalous in standardized
units.  $A_1$ represents the moisture pathway through kernel-integrated
relative humidity, and $A_2$ represents the instability pathway through the
imbalance between kernel-integrated equivalent potential temperature and its
saturation value.  The cubic exponent amplifies the dominant pathway, and
$s_y\,a$ sets the overall sensitivity.  The offset $d_0$ in $A_2$ absorbs
the training-set climatology of the instability measure and the fitted
constant $c$.

Unlike SR-BL, which reduces to three interpretable constants, SR-ATM operates
on kernel-integrated features that combine information across multiple
vertical levels. Its physical form retains the standardization statistics as
explicit parameters rather than absorbing them into a single threshold
and sensitivity.

\vspace{0.3em}
\noindent\textit{Numerical values.}  With the optimized constants from
Table~S2 ($a = \text{[TBD]}$, $b = \text{[TBD]}$, $c = \text{[TBD]}$),
kernel parameters from Table~S1, and training-set statistics:
%
\begin{equation*}
  s_y\,a = \text{[TBD]}\,,\qquad
  d_0 = \text{[TBD]}\,.
\end{equation*}

% -------------------------------------------------------------------------

\vspace{1.0em}
\textbf{SR-SFC: Surface Flux Correction.}
SR-SFC extends SR-ATM by adding a surface flux correction term.  During
symbolic regression, the optimized SR-ATM prediction $z_{\mathrm{ATM}}$
enters as a fixed input, so that only the correction is learned.  The
discovered equation in standardized space (Table~S2) is
%
\begin{equation}
  f(\mathbf{x}) \;=\; z_{\mathrm{ATM}}
    + a\,\widetilde{\mathrm{SHF}}\,(b - \mathrm{LF})
    + c\,\widetilde{\mathrm{LHF}}\,,
  \label{eq:s-srsfc-std}
\end{equation}
%
with optimized constants $a$, $b$, and $c$ (Table~SX), where
$z_{\mathrm{ATM}} = z_{\min} + \max(f_{\mathrm{ATM}}(\mathbf{x}),\,0)$ is
the SR-ATM prediction held fixed during optimization.

The surface flux predictors are scalar quantities standardized as
$\widetilde{\mathrm{SHF}} = (\mathrm{SHF} -
\mu_{\mathrm{SHF}})/s_{\mathrm{SHF}}$ and $\widetilde{\mathrm{LHF}} =
(\mathrm{LHF} - \mu_{\mathrm{LHF}})/s_{\mathrm{LHF}}$.  LF is not
standardized.  Substituting:
%
\begin{align}
  f(\mathbf{x}) &= z_{\mathrm{ATM}}
    + \frac{a}{s_{\mathrm{SHF}}}(\mathrm{SHF} - \mu_{\mathrm{SHF}})
      (b - \mathrm{LF})
    + \frac{c}{s_{\mathrm{LHF}}}(\mathrm{LHF} - \mu_{\mathrm{LHF}})
    \notag\\[4pt]
    &= z_{\mathrm{ATM}}
    + \frac{a}{s_{\mathrm{SHF}}}\,\mathrm{SHF}\,(b - \mathrm{LF})
    - \frac{a\,\mu_{\mathrm{SHF}}}{s_{\mathrm{SHF}}}(b - \mathrm{LF})
    + \frac{c}{s_{\mathrm{LHF}}}\,\mathrm{LHF}
    - \frac{c\,\mu_{\mathrm{LHF}}}{s_{\mathrm{LHF}}}\,.
  \label{eq:s-srsfc-expand}
\end{align}
%
Define grouped physical constants:
%
\begin{equation}
  \begin{aligned}
    \alpha_S &= a\,/\,s_{\mathrm{SHF}} &&\text{(SHF sensitivity)}\,,\\
    \alpha_L &= c\,/\,s_{\mathrm{LHF}} &&\text{(LHF sensitivity)}\,,\\
    \delta_S &= a\,\mu_{\mathrm{SHF}}\,/\,s_{\mathrm{SHF}}
             &&\text{(SHF climatological offset)}\,,\\
    \delta_L &= c\,\mu_{\mathrm{LHF}}\,/\,s_{\mathrm{LHF}}
             &&\text{(LHF climatological offset)}\,.
  \end{aligned}
  \label{eq:s-srsfc-constants}
\end{equation}
%
The full prediction is then
%
\begin{equation}
  \boxed{P \;=\; \max\!\Bigl(\exp\!\bigl(
    s_y\bigl[f_{\mathrm{ATM}}
    + \alpha_S\,\mathrm{SHF}\,(b - \mathrm{LF})
    + \alpha_L\,\mathrm{LHF}
    - \delta_S\,(b - \mathrm{LF})
    - \delta_L\bigr]\bigr) - 1,\;\;0\Bigr),}
  \label{eq:s-srsfc-phys}
\end{equation}
%
where $f_{\mathrm{ATM}} = \max(f_{\mathrm{ATM}}(\mathbf{x}),\,0)$ and the
inner $\max(\cdot,\,0)$ on $f_{\mathrm{ATM}}$ and the outer
$\max(\cdot,\,0)$ on $P$ arise from the shared prediction structure
(Equation~\ref{eq:s-zfloor}).  Note that $z_{\mathrm{ATM}}$ in the
standardized equation already includes $z_{\min}$, so the $\mu_y$
cancellation described above carries through to the composite prediction.

The correction term has a clear physical interpretation.  The SHF contribution
is modulated by $(b - \mathrm{LF})$, where $b \approx 0.74$ (Table~S2)
implies that SHF enhances precipitation primarily over ocean ($\mathrm{LF}
\approx 0$) and its influence vanishes near $\mathrm{LF} \approx b$.  The
LHF contribution enters as a linear additive term with no land--ocean
modulation.

\vspace{0.3em}
\noindent\textit{Numerical values.}  With the optimized constants from
Table~S2 ($a = \text{[TBD]}$, $b = \text{[TBD]}$, $c = \text{[TBD]}$)
and training-set statistics:
%
\begin{equation*}
  \alpha_S = \text{[TBD]}\,,\qquad
  \alpha_L = \text{[TBD]}\,,\qquad
  \delta_S = \text{[TBD]}\,,\qquad
  \delta_L = \text{[TBD]}\,.
\end{equation*}

% -------------------------------------------------------------------------

\vspace{1.0em}
\textbf{SR-ALL: Combined Atmospheric and Surface Equation.}
SR-ALL extends SR-ATM by incorporating surface fluxes directly into the
atmospheric equation rather than as a separate additive correction.  Like
SR-SFC, $z_{\mathrm{ATM}}$ enters as a fixed input.  The discovered equation
in standardized space (Table~S2) is
%
\begin{equation}
  f(\mathbf{x}) \;=\; z_{\mathrm{ATM}}
    + (\hat{\theta}_e + a\,\widetilde{\mathrm{SHF}})\,(b - \mathrm{LF})^3
    + c\,,
  \label{eq:s-srall-std}
\end{equation}
%
with optimized constants $a$, $b$, and $c$ (Table~SX).

Substituting the standardized forms $\hat{\theta}_e = (\bar{\theta}_e -
\mu_{\theta_e})/s_{\theta_e}$ and $\widetilde{\mathrm{SHF}} = (\mathrm{SHF}
- \mu_{\mathrm{SHF}})/s_{\mathrm{SHF}}$:
%
\begin{align}
  f(\mathbf{x}) &= z_{\mathrm{ATM}}
    + \biggl(\frac{\bar{\theta}_e - \mu_{\theta_e}}{s_{\theta_e}}
    + a\,\frac{\mathrm{SHF} - \mu_{\mathrm{SHF}}}{s_{\mathrm{SHF}}}\biggr)
    (b - \mathrm{LF})^3 + c \notag\\[4pt]
    &= z_{\mathrm{ATM}}
    + \biggl(\frac{\bar{\theta}_e}{s_{\theta_e}}
    + \frac{a\,\mathrm{SHF}}{s_{\mathrm{SHF}}}
    - \underbrace{\frac{\mu_{\theta_e}}{s_{\theta_e}}
    - \frac{a\,\mu_{\mathrm{SHF}}}{s_{\mathrm{SHF}}}}_{\displaystyle
    \equiv\; e_0}\biggr)(b - \mathrm{LF})^3 + c\,.
  \label{eq:s-srall-expand}
\end{align}
%
Define grouped physical constants:
%
\begin{equation}
  \begin{aligned}
    \gamma_\theta &= 1\,/\,s_{\theta_e}
      &&\text{($\bar{\theta}_e$ sensitivity per unit)}\,,\\
    \gamma_S &= a\,/\,s_{\mathrm{SHF}}
      &&\text{(SHF sensitivity per unit)}\,,\\
    e_0 &= \mu_{\theta_e}\,/\,s_{\theta_e}
         + a\,\mu_{\mathrm{SHF}}\,/\,s_{\mathrm{SHF}}
      &&\text{(climatological offset)}\,.
  \end{aligned}
  \label{eq:s-srall-constants}
\end{equation}
%
The full prediction is then
%
\begin{equation}
  \boxed{P \;=\; \max\!\Bigl(\exp\!\bigl(
    s_y\bigl[f_{\mathrm{ATM}}
    + \bigl(\gamma_\theta\,\bar{\theta}_e
    + \gamma_S\,\mathrm{SHF} - e_0\bigr)
    (b - \mathrm{LF})^3 + c\bigr]\bigr) - 1,\;\;0\Bigr),}
  \label{eq:s-srall-phys}
\end{equation}
%
where $f_{\mathrm{ATM}}$ is defined as in SR-SFC above.

In SR-ALL, the correction couples the atmospheric state ($\bar{\theta}_e$)
with the surface energy input (SHF) multiplicatively, scaled by a cubic
function of land fraction.  The $(b - \mathrm{LF})^3$ factor with $b \approx
0.73$ (Table~S2) produces a sign change at $\mathrm{LF} \approx b$, so that
the correction has opposite sign over ocean ($\mathrm{LF} < b$) and over
land ($\mathrm{LF} > b$).  The cubic dependence on LF amplifies this contrast
relative to the linear modulation in SR-SFC.

\vspace{0.3em}
\noindent\textit{Numerical values.}  With the optimized constants from
Table~S2 ($a = \text{[TBD]}$, $b = \text{[TBD]}$, $c = \text{[TBD]}$)
and training-set statistics:
%
\begin{equation*}
  \gamma_\theta = \text{[TBD]}\,,\qquad
  \gamma_S = \text{[TBD]}\,,\qquad
  e_0 = \text{[TBD]}\,.
\end{equation*}

% -------------------------------------------------------------------------

\vspace{1.0em}
\textbf{Summary.}
Table~SX collects the physical constants for all four symbolic regression
equations.  SR-BL admits a clean three-parameter form (Equation~\ref{eq:s-srbl-phys})
because its single scalar predictor requires only one standardization.  SR-ATM,
SR-SFC, and SR-ALL retain the standardization statistics as explicit grouped
constants because they operate on multiple predictors with different units and
scales.  In all cases, the physical forms are exact algebraic rearrangements of
the standardized equations in Table~S2 and require no retraining or
approximation.

% -------------------------------------------------------------------------
% Placeholder for a summary table of physical constants.
% Uncomment and populate once results are finalized.
% -------------------------------------------------------------------------
%
% \begin{table}[h]
% \centering
% \caption{Physical constants for each symbolic regression equation.
%   All values are derived algebraically from the optimized standardized
%   constants (Table~S2) and the training-set statistics.}
% \label{tab:s-physical-constants}
% \begin{tabular}{llrl}
% \hline
% Equation & Constant & Value & Interpretation \\
% \hline
% SR-BL & $B_c$    & [TBD] m\,s$^{-2}$ & Critical buoyancy threshold \\
%        & $\alpha$ & [TBD] (m\,s$^{-2}$)$^{-3}$ & Exponential sensitivity \\
%        & $\beta$  & [TBD] & Log-space offset \\
% \hline
% SR-ATM & $s_y a$  & [TBD] & Overall sensitivity \\
%        & $d_0$    & [TBD] & Instability offset \\
% \hline
% SR-SFC & $\alpha_S$ & [TBD] & SHF sensitivity \\
%        & $\alpha_L$ & [TBD] & LHF sensitivity \\
%        & $\delta_S$ & [TBD] & SHF offset \\
%        & $\delta_L$ & [TBD] & LHF offset \\
% \hline
% SR-ALL & $\gamma_\theta$ & [TBD] & $\bar{\theta}_e$ sensitivity \\
%        & $\gamma_S$      & [TBD] & SHF sensitivity \\
%        & $e_0$           & [TBD] & Climatological offset \\
% \hline
% \end{tabular}
% \end{table}
